% Canonical split-Cartan SU(r) extension module, 2026-09-03.
% This is a delta to cusp_native_full_flux_magnetic.tex.  It imports that
% module's native torus, full-flux line bundle, orientations, exact magnetic
% frequencies, and SU(2) calculation.  It deliberately does not duplicate
% the radial direct-integral operator or the charge-two SU(2) proof.

\section{The native full-flux split-Cartan \texorpdfstring{\(SU(r)\)}{SU(r)} extension}
\label{cusp-native-full-flux-suN-section}

This section fixes an integer \(r\geq2\) and extends the full-flux
construction of \cref{cusp-native-full-flux-section} from its already proved
split \(SU(2)\) case to arbitrary split-Cartan \(SU(r)\) backgrounds.  Every
object remains on the same regular flat four-torus
\(F_T=\mathbb R^4/P_T\mathbb Z^4\), with the unrescaled Euclidean quotient
metric.  In particular, the integral two-form \(\eta_T\), the line bundle
\(\mathscr L_{\eta,T}\), the two orientations, and the exact quantities
\[
 \begin{gathered}
  L_T,\quad q_T,\quad m_T,\quad D_T=L_T q_T+6m_T^2>0,\\
  S_T=6L_T-q_T,\qquad
  \Xi_T=\sqrt{(6L_T+q_T)^2+144m_T^2}.
 \end{gathered}
\]
are precisely those of
\cref{cusp-native-full-flux-canonical-scalars,%
cusp-native-full-flux-S-Xi}.  The charge-one frequencies are
\begin{equation}
 \mathcal B_{T,+}=\frac{\pi(S_T+\Xi_T)}{D_T},
 \qquad
 \mathcal B_{T,-}=\frac{\pi(\Xi_T-S_T)}{D_T},
 \label{cusp-native-full-flux-suN-signed-frequencies}
\end{equation}
and we shall use the globally ordered notation
\begin{equation}
 \begin{aligned}
  \mathcal B_{T,\max}
  &=\max\{\mathcal B_{T,+},\mathcal B_{T,-}\},&
  \mathcal B_{T,\min}
  &=\min\{\mathcal B_{T,+},\mathcal B_{T,-}\},\\
  \rho_T^*&=\frac{\mathcal B_{T,\max}}{\mathcal B_{T,\min}}.
 \end{aligned}
 \label{cusp-native-full-flux-suN-ordered-frequencies}
\end{equation}
Thus every assertion below remains valid on either side of \(S_T=0\).

\subsection{The split bundle, determinant, and Chern number}

Fix a nonzero trace-zero integral weight vector
\begin{equation}
 \mathbf m=(m_1,\ldots,m_r)\in\mathbb Z^r,
 \qquad \sum_{j=1}^r m_j=0,
 \qquad s_2(\mathbf m):=\sum_{j=1}^r m_j^2.
 \label{cusp-native-full-flux-suN-weight-vector}
\end{equation}
Define
\begin{equation}
 \mathscr E_{\mathbf m,T}:=
 \bigoplus_{j=1}^r\mathscr L_{\eta,T}^{\otimes m_j},
 \qquad
 H_{\mathbf m}:=i\operatorname{diag}(m_1,\ldots,m_r)
 \in\mathfrak{su}(r).
 \label{cusp-native-full-flux-suN-split-bundle}
\end{equation}

\begin{proposition}[Parallel determinant trivialization and exact \(c_2\)]
\label{cusp-native-full-flux-suN-topology}
The ordered splitting in
\cref{cusp-native-full-flux-suN-split-bundle} has a canonical parallel
determinant trivialization.  With \(c=[\eta_T]\),
\begin{equation}
 c_1(\mathscr E_{\mathbf m,T})=0,
 \qquad
 c_2(\mathscr E_{\mathbf m,T})
 =-\frac{s_2(\mathbf m)}2c^2.
 \label{cusp-native-full-flux-suN-Chern-classes}
\end{equation}
The marked and physical-complex Chern numbers are respectively
\begin{equation}
 \left\langle c_2,[F_T]_{\mathrm{mark}}\right\rangle=-6s_2(\mathbf m),
 \qquad
 k_{\mathbf m}:=
 \left\langle c_2,[F_T]_{x\text{-}\mathrm{complex}}\right\rangle
 =6s_2(\mathbf m)>0.
 \label{cusp-native-full-flux-suN-Chern-numbers}
\end{equation}
Moreover \(s_2(\mathbf m)\) is even.  Hence this direct-sum family realizes
only \(k_{\mathbf m}\in12\mathbb Z_{>0}\), not all topological \(SU(r)\)
sectors.
\end{proposition}

\begin{proof}
Let \(j_n^\eta(y)\) denote the factor of automorphy in
\cref{cusp-native-full-flux-cocycle}.  The transition matrix of the ordered
sum is
\[
 \operatorname{diag}\bigl((j_n^\eta)^{m_1},\ldots,
 (j_n^\eta)^{m_r}\bigr).
\]
Its determinant is \((j_n^\eta)^{\sum_j m_j}=1\) identically.  The tensor
product of the ordered unit frames therefore defines a global determinant
section, and the determinant connection has connection form proportional to
\(\sum_j m_j=0\); the section is parallel.

The splitting principle gives
\[
 c(\mathscr E_{\mathbf m,T})=\prod_{j=1}^r(1+m_j c),
 \qquad
 c_2=\sum_{i<j}m_i m_j c^2.
\]
Since
\[
 2\sum_{i<j}m_i m_j
 =\left(\sum_j m_j\right)^2-\sum_j m_j^2=-s_2(\mathbf m),
\]
this proves \cref{cusp-native-full-flux-suN-Chern-classes}.  The two
evaluations follow from
\cref{cusp-native-full-flux-two-orientation-pairings}; retaining its
orientation reversal is essential for the sign.  Finally
\[
 s_2(\mathbf m)=\sum_j m_j^2\equiv\sum_j m_j\equiv0\pmod2,
\]
because \(a^2\equiv a\pmod2\).  This proves the divisibility statement.
\end{proof}

The diagonal connection induced by the summands is
\begin{equation}
 \nabla^{\mathbf m,T}
 =\mathrm d-2\pi(6y^1\,\mathrm dy^4+y^2\,\mathrm dy^3)H_{\mathbf m},
 \qquad
 F_{\mathbf m,T}=-2\pi H_{\mathbf m}\eta_T.
 \label{cusp-native-full-flux-suN-connection-curvature}
\end{equation}
We retain the invariant inner product and functional normalization
\begin{equation}
 \begin{gathered}
  \langle X,Y\rangle_{\mathfrak{su}(r)}
  =-\frac12\operatorname{Tr}(XY),
  \qquad
  \lVert H_{\mathbf m}\rVert^2=\frac12s_2(\mathbf m),\\
  \operatorname{YM}_T(A)=\frac1{2g_{\mathrm{YM}}^2}
  \int_{F_T}|F_A|^2\,\mathrm d\operatorname{vol}_T.
 \end{gathered}
 \label{cusp-native-full-flux-suN-normalizations}
\end{equation}

\begin{theorem}[Yang--Mills criticality and exact action]
\label{cusp-native-full-flux-suN-action-theorem}
The connection \(\nabla^{\mathbf m,T}\) is non-flat, reducible, and
Yang--Mills.  Its action is
\begin{align}
 \operatorname{YM}_T(\nabla^{\mathbf m,T})
 &=\frac{\pi^2s_2(\mathbf m)}{g_{\mathrm{YM}}^2D_T}
   \bigl(q_T^2+36L_T^2+72m_T^2\bigr)
 \label{cusp-native-full-flux-suN-action-coordinate}\\
 &=\frac{2\pi^2k_{\mathbf m}}{g_{\mathrm{YM}}^2}
   +\frac{\pi^2s_2(\mathbf m)S_T^2}
          {g_{\mathrm{YM}}^2D_T}.
 \label{cusp-native-full-flux-suN-action-excess}
\end{align}
On the canonical cusp tail,
\begin{equation}
 \operatorname{YM}_T(\nabla^{\mathbf m,T})
 =\frac{\pi^2s_2(\mathbf m)}{g_{\mathrm{YM}}^2}
 \left(37+\frac{35B}{T}+O(T^{-2})\right).
 \label{cusp-native-full-flux-suN-action-asymptotic}
\end{equation}
\end{theorem}

\begin{proof}
Both \(H_{\mathbf m}\) and \(\eta_T\) are parallel for the diagonal
connection and flat quotient metric, and their brackets vanish.  Hence
\(d_{\nabla^{\mathbf m,T}}^*F_{\mathbf m,T}=0\).  Non-flatness follows
from \(\mathbf m\ne0\) and \(\eta_T\ne0\).

In the physical orthonormal frame used in
\cref{cusp-native-full-flux-normal-form},
\[
 |\eta_T|^2
 =\frac{\mathcal B_{T,+}^2+\mathcal B_{T,-}^2}{4\pi^2}.
\]
Using \cref{cusp-native-full-flux-suN-normalizations},
\(\operatorname{vol}(F_T)=D_T\), and the exact frequency square sum from
the native module yields
\cref{cusp-native-full-flux-suN-action-coordinate}.  The identity
\[
 q_T^2+36L_T^2+72m_T^2
 =12(L_T q_T+6m_T^2)+(6L_T-q_T)^2
 =12D_T+S_T^2
\]
and \(k_{\mathbf m}=6s_2(\mathbf m)\) give
\cref{cusp-native-full-flux-suN-action-excess}.  Substituting the retained
expansions of \(L_T,q_T,m_T,D_T\), without suppressing \(H,B,m\), gives
\[
 \frac{q_T^2+36L_T^2+72m_T^2}{D_T}
 =37+\frac{35B}{T}+O(T^{-2}),
\]
which proves the last formula.
\end{proof}

\subsection{Every nonzero integral root charge}

The native module proves the charge-one ladder decomposition and the
charge-two specialization used by its \(SU(2)\) theorem.  The following
lemma records, and proves, the exact integer-charge extension needed here.

\begin{lemma}[Charge-\(d\) Landau decomposition]
\label{cusp-native-full-flux-suN-charge-d-Landau}
Let \(d\in\mathbb Z\setminus\{0\}\), put \(q=|d|\), and equip
\(\mathscr L_{\eta,T}^{\otimes d}\) with its tensor-power connection.  Its
scalar covariant Laplacian has joint levels
\begin{equation}
 q\bigl((2a+1)\mathcal B_{T,+}+(2b+1)\mathcal B_{T,-}\bigr),
 \qquad a,b\in\mathbb Z_{\ge0},
 \label{cusp-native-full-flux-suN-charge-d-levels}
\end{equation}
and every joint space has complex dimension \(6q^2\).  If different pairs
\((a,b)\) give the same numerical value, their orthogonal joint spaces are
to be summed.
\end{lemma}

\begin{proof}
The curvature is \(-2\pi i d\eta_T\).  In the polar construction of
\cref{cusp-native-full-flux-landau-theorem}, replace the auxiliary complex
structure by
\(J_{d,T}=\operatorname{sgn}(d)J_{\eta,T}\).  Then \(d\eta_T\) is positive
of type \((1,1)\), so the vacuum is the holomorphic-section space of the
positive bundle \(\mathscr L_{\eta,T}^{\otimes d}\) on that auxiliary
complex torus.  Its tangent and canonical bundles are trivial.  Kodaira
vanishing and Hirzebruch--Riemann--Roch give
\[
 \dim_{\mathbb C}H^0
 =\frac12\int_{[F_T]_{d\eta\text{-}\mathrm{symp}}}(d\eta_T)^2
 =\frac12q^2\int_{[F_T]_{\mathrm{mark}}}\eta_T^2
 =6q^2.
\]
Here changing \(d\) to \(-d\) does not change the four-dimensional
symplectic orientation.  Scaling the curvature by \(d\) scales both
positive oscillator frequencies by \(q\), while the two canonical
commutation relations and the lowering argument in
\cref{cusp-native-full-flux-landau-theorem} are unchanged after this
rescaling.  Thus every joint level is unitarily obtained from the vacuum and
has the asserted energy and dimension.  Orthogonal joint spaces remain a
direct sum when their numerical energies coincide.
\end{proof}

\subsection{Closed gauge-fixed and Coulomb operators}

Write \(A=\nabla^{\mathbf m,T}\) in this subsection.  Let
\begin{equation}
 \mathscr H_{\mathbf m,T}^1
 =L^2\Omega^1(F_T;\operatorname{ad}\mathscr E_{\mathbf m,T}),
 \qquad
 \mathscr H_{\mathbf m,T}^0
 =L^2\Omega^0(F_T;\operatorname{ad}\mathscr E_{\mathbf m,T})
 \label{cusp-native-full-flux-suN-real-Hilbert-spaces}
\end{equation}
as real Hilbert spaces.  Define the maximal distributional codifferential
by
\begin{align}
 \operatorname{Dom}D_{\mathbf m,T}^{\max}
 =\{a\in\mathscr H_{\mathbf m,T}^1:\;&
 \exists u\in\mathscr H_{\mathbf m,T}^0\text{ such that}\notag\\[-1mm]
 &\langle a,d_A\phi\rangle_{L^2}
 =\langle u,\phi\rangle_{L^2}
 \quad\text{for every }\notag\\[-1mm]
 &\phi\in C^\infty(F_T;\operatorname{ad}\mathscr E_{\mathbf m,T})\},
 \label{cusp-native-full-flux-suN-max-codiff-domain}\\
 D_{\mathbf m,T}^{\max}a={}&u,
 \qquad
 \mathscr S_{\mathbf m,T}:=\ker D_{\mathbf m,T}^{\max}.
 \label{cusp-native-full-flux-suN-Coulomb-slice}
\end{align}
Density of smooth test scalars makes \(u\) unique.  If
\(a_j\to a\) and \(D_{\mathbf m,T}^{\max}a_j\to u\) in \(L^2\), passage to
the limit in the defining identity gives
\(D_{\mathbf m,T}^{\max}a=u\).  Thus \(D_{\mathbf m,T}^{\max}\) is closed
and \(\mathscr S_{\mathbf m,T}\) is a closed subspace.

Let \(\Delta_A^{(0)}\) be the nonnegative self-adjoint scalar operator of
the closed form
\begin{equation}
 \mathfrak q_{\mathbf m,T}^{(0)}[\phi,\psi]
 =\langle d_A\phi,d_A\psi\rangle_{L^2},
 \qquad
 \operatorname{Dom}\mathfrak q_{\mathbf m,T}^{(0)}=H^1\Omega^0.
 \label{cusp-native-full-flux-suN-scalar-form}
\end{equation}
It has domain \(H^2\Omega^0\) and equals \(d_A^*d_A\).  In a constant
orthonormal frame set
\begin{equation}
 (V_{\mathbf m,T}a)_\nu
 =-2\sum_{\mu=1}^4[F_{\mathbf m,T,\nu\mu},a_\mu].
 \label{cusp-native-full-flux-suN-spin-endomorphism}
\end{equation}
Invariance of the compact Lie-algebra inner product and
\(F_{\nu\mu}=-F_{\mu\nu}\) show directly that \(V_{\mathbf m,T}\) is a
bounded symmetric endomorphism.  Define
\begin{equation}
 \mathfrak j_{\mathbf m,T}^{\mathrm{gf}}[a,b]
 =\langle\nabla_A a,\nabla_A b\rangle_{L^2}
  +\langle V_{\mathbf m,T}a,b\rangle_{L^2},
 \qquad
 \operatorname{Dom}\mathfrak j_{\mathbf m,T}^{\mathrm{gf}}
 =H^1\Omega^1.
 \label{cusp-native-full-flux-suN-Jacobi-form}
\end{equation}

\begin{proposition}[Closed realization and exact Coulomb reduction]
\label{cusp-native-full-flux-suN-Coulomb-reduction}
The form in \cref{cusp-native-full-flux-suN-Jacobi-form} is densely defined,
closed, symmetric, and bounded below.  Its self-adjoint operator has compact
resolvent, domain \(H^2\Omega^1\), and expression
\begin{equation}
 (\mathcal J_{\mathbf m,T}^{\mathrm{gf}}a)_\nu
 =-\sum_{\mu=1}^4\nabla_{A,e_\mu}\nabla_{A,e_\mu}a_\nu
  -2\sum_{\mu=1}^4[F_{\mathbf m,T,\nu\mu},a_\mu].
 \label{cusp-native-full-flux-suN-Jacobi-expression}
\end{equation}
For every \(z\in\mathbb C\setminus\mathbb R\) and
\(f\in\operatorname{Dom}D_{\mathbf m,T}^{\max}\),
\begin{equation}
 D_{\mathbf m,T}^{\max}
 (\mathcal J_{\mathbf m,T}^{\mathrm{gf}}-z)^{-1}f
 =(\Delta_A^{(0)}-z)^{-1}D_{\mathbf m,T}^{\max}f.
 \label{cusp-native-full-flux-suN-resolvent-intertwining}
\end{equation}
Consequently \(\mathscr S_{\mathbf m,T}\) reduces
\(\mathcal J_{\mathbf m,T}^{\mathrm{gf}}\), and
\begin{equation}
 \mathcal J_{\mathbf m,T}
 :=\mathcal J_{\mathbf m,T}^{\mathrm{gf}}
   \big|_{\mathscr S_{\mathbf m,T}},
 \qquad
 \operatorname{Dom}\mathcal J_{\mathbf m,T}
 =H^2\Omega^1\cap\mathscr S_{\mathbf m,T}
 \label{cusp-native-full-flux-suN-Coulomb-operator}
\end{equation}
is self-adjoint, with compact resolvent and closed form
\(\mathfrak j_{\mathbf m,T}\) equal to the restriction of
\(\mathfrak j_{\mathbf m,T}^{\mathrm{gf}}\) to
\(H^1\Omega^1\cap\mathscr S_{\mathbf m,T}\).  The physical Yang--Mills
Hessian on this slice is
\begin{equation}
 \operatorname{YM}_T''\big|_A
 =g_{\mathrm{YM}}^{-2}\mathfrak j_{\mathbf m,T}.
 \label{cusp-native-full-flux-suN-physical-Hessian}
\end{equation}
Thus an eigenvalue \(\nu\) of \(\mathcal J_{\mathbf m,T}\) is the physical
Hessian eigenvalue \(\nu/g_{\mathrm{YM}}^2\), while its sign and Morse-index
contribution are unchanged.
\end{proposition}

\begin{proof}
The first term of
\cref{cusp-native-full-flux-suN-Jacobi-form}, after addition of the
\(L^2\)-norm, is an equivalent \(H^1\)-norm.  The second term is a bounded
symmetric form.  The bounded perturbation theorem for closed forms and the
representation theorem therefore give a lower-bounded self-adjoint
operator.  Its principal symbol is \(|\xi|^2\) times the identity, so
elliptic regularity gives the \(H^2\) domain.  Rellich compactness on the
compact torus gives compact resolvent.

The two-parameter curvature expansion and the flat Weitzenb\"ock identity,
with the inner product and functional normalization in
\cref{cusp-native-full-flux-suN-normalizations}, give
\[
 \operatorname{YM}_T''\big|_A(a,b)
 +g_{\mathrm{YM}}^{-2}\langle d_A^*a,d_A^*b\rangle
 =g_{\mathrm{YM}}^{-2}
   \mathfrak j_{\mathbf m,T}^{\mathrm{gf}}[a,b].
\]
This proves the Hessian normalization after restriction to Coulomb gauge.

It remains to prove that the restriction is an operator rather than a formal
constraint.  For a smooth adjoint scalar \(\phi\), write
\(\nabla_\mu=\nabla_{A,e_\mu}\).  Commuting all three derivatives and using
the Yang--Mills equation gives
\begin{align*}
 &(\mathcal J_{\mathbf m,T}^{\mathrm{gf}}d_A\phi
       -d_A\Delta_A^{(0)}\phi)_\nu\\
 &\quad=-\nabla_\mu\nabla_\mu\nabla_\nu\phi
       +\nabla_\nu\nabla_\mu\nabla_\mu\phi
       -2[F_{\nu\mu},\nabla_\mu\phi]\\
 &\quad=-[\nabla_\mu F_{\mu\nu},\phi]
       =[(d_A^*F_A)_\nu,\phi]=0.
\end{align*}
Formal adjunction yields, on smooth one-forms,
\begin{equation}
 \mathcal J_{\mathbf m,T}^{\mathrm{gf}}d_A
 =d_A\Delta_A^{(0)},
 \qquad
 d_A^*\mathcal J_{\mathbf m,T}^{\mathrm{gf}}
 =\Delta_A^{(0)}d_A^*.
 \label{cusp-native-full-flux-suN-formal-intertwining}
\end{equation}
For
\(u=(\mathcal J_{\mathbf m,T}^{\mathrm{gf}}-z)^{-1}f\), the operator domain
gives \(u\in H^2\), hence \(D_{\mathbf m,T}^{\max}u\in H^1\).  Applying the
second identity distributionally to
\((\mathcal J_{\mathbf m,T}^{\mathrm{gf}}-z)u=f\) gives
\[
 (\Delta_A^{(0)}-z)D_{\mathbf m,T}^{\max}u
 =D_{\mathbf m,T}^{\max}f.
\]
Scalar elliptic regularity and the scalar resolvent give
\cref{cusp-native-full-flux-suN-resolvent-intertwining}.  The two conjugate
resolvents preserve \(\mathscr S_{\mathbf m,T}\); adjunction then makes its
orthogonal complement invariant as well.  The slice is reducing, so the
self-adjoint part has the operator and form domains displayed above.
\end{proof}

\subsection{Rootwise spectrum, coincidences, index, and nullity}

For \(i\ne j\), the matrix unit \(E_{ij}\) obeys
\[
 [H_{\mathbf m},E_{ij}]=i(m_i-m_j)E_{ij},
\]
so its complex root line is
\begin{equation}
 \operatorname{Hom}(\mathscr L_{\eta,T}^{m_j},
                     \mathscr L_{\eta,T}^{m_i})
 =\mathscr L_{\eta,T}^{m_i-m_j},
 \qquad q_{ij}:=|m_i-m_j|.
 \label{cusp-native-full-flux-suN-root-charge}
\end{equation}
The factor \(q_{ij}\) is the complete charge normalization; for the existing
\(SU(2)\) weights \((1,-1)\), it gives \(q_{12}=2\).

Let the distinct numerical values among \(m_1,\ldots,m_r\) have
multiplicities \(r_1,\ldots,r_\ell\).  The neutral centralizer is
\begin{equation}
 \mathfrak z_{\mathbf m}
 =\{X\in\mathfrak{su}(r):[H_{\mathbf m},X]=0\}
 \cong\mathfrak s\!\left(\bigoplus_{a=1}^{\ell}\mathfrak u(r_a)\right),
 \qquad
 \dim_{\mathbb R}\mathfrak z_{\mathbf m}
 =\sum_{a=1}^{\ell}r_a^2-1.
 \label{cusp-native-full-flux-suN-centralizer}
\end{equation}

\begin{theorem}[Complete negative spectrum and numerical multiplicities]
\label{cusp-native-full-flux-suN-Morse-theorem}
Put
\begin{equation}
 N_-(T):=\#\{n\in\mathbb Z_{\ge0}:2n+1<\rho_T^*\}
 =\left\lceil\frac{\rho_T^*-1}{2}\right\rceil.
 \label{cusp-native-full-flux-suN-N-minus}
\end{equation}
For each positive root \(i<j\) with \(q_{ij}>0\), the complete negative
spectrum of the normalized Coulomb operator is
\begin{equation}
 \nu_{ij,n}(T)
 =q_{ij}\bigl((2n+1)\mathcal B_{T,\min}
                  -\mathcal B_{T,\max}\bigr),
 \qquad 0\le n<N_-(T).
 \label{cusp-native-full-flux-suN-negative-bands}
\end{equation}
The contribution of that positive-root pair to each displayed band has real
dimension \(12q_{ij}^2\).  For an actual numerical eigenvalue
\(\lambda<0\), its full real multiplicity is therefore
\begin{equation}
 \operatorname{mult}_{\mathbb R}(\lambda)
 =\sum_{\substack{i<j,\ q_{ij}>0,\ 0\le n<N_-(T)\\
                   \nu_{ij,n}(T)=\lambda}}
   12q_{ij}^2.
 \label{cusp-native-full-flux-suN-coincident-multiplicity}
\end{equation}
In particular, coincident root charges or accidental coincidences between
different \(q_{ij},n\) add rather than being counted as one band.  The real
Morse index is
\begin{equation}
 \boxed{
 \operatorname{ind}_{\mathbb R}\mathcal J_{\mathbf m,T}
 =12N_-(T)\sum_{i<j}(m_i-m_j)^2
 =12rN_-(T)s_2(\mathbf m).}
 \label{cusp-native-full-flux-suN-index}
\end{equation}
Every eigenvalue in
\cref{cusp-native-full-flux-suN-negative-bands} is divided by
\(g_{\mathrm{YM}}^2\) for the physical Hessian.
\end{theorem}

\begin{proof}
On the \(E_{ij}\) line, the scalar oscillator is the charge-\(q_{ij}\)
operator of \cref{cusp-native-full-flux-suN-charge-d-Landau}.  In a
skew-normal orthonormal coframe, the four circular one-form polarizations
diagonalize \cref{cusp-native-full-flux-suN-spin-endomorphism}; the two
polarizations in a plane of charge-\(q_{ij}\) frequency
\(q_{ij}\mathcal B\) shift the scalar energy by
\(\pm2q_{ij}\mathcal B\).  The unfavorable polarizations are positive.
For the favorable polarization in the high-frequency plane, high-plane
oscillator number zero leaves
\[
 q_{ij}\bigl((2n+1)\mathcal B_{T,\min}
             -\mathcal B_{T,\max}\bigr).
\]
Increasing its high-plane oscillator number makes it positive.  The
favorable low-frequency polarization has bottom
\(q_{ij}(\mathcal B_{T,\max}-\mathcal B_{T,\min})\ge0\), and increasing
either number cannot lower it.  This proves that the list is exhaustive.

Every joint scalar space on one complex root line has dimension
\(6q_{ij}^2\).  The opposite root is its anti-linear conjugate.  The pair is
the complexification of a real root subspace, so its real dimension is
\(12q_{ij}^2\).  Root subspaces and joint oscillator spaces are orthogonal;
hence equality of numerical eigenvalues produces precisely the direct sum
in \cref{cusp-native-full-flux-suN-coincident-multiplicity}.

The strict inequality in
\cref{cusp-native-full-flux-suN-N-minus} counts the integers correctly even
when \(\rho_T^*\) is an odd integer.  Summing dimensions gives the first
index formula.  Finally,
\[
 \sum_{i<j}(m_i-m_j)^2
 =r\sum_i m_i^2-\left(\sum_i m_i\right)^2
 =rs_2(\mathbf m),
\]
which gives the second formula.

It remains to justify Coulomb gauge.  If
\(\mathcal J_{\mathbf m,T}^{\mathrm{gf}}a=\lambda a\) with
\(\lambda<0\), elliptic regularity and
\cref{cusp-native-full-flux-suN-formal-intertwining} give
\[
 \Delta_A^{(0)}d_A^*a=\lambda d_A^*a.
\]
The scalar operator is nonnegative, so \(d_A^*a=0\).  Thus the negative
gauge-fixed and Coulomb eigenspaces coincide, completing the proof.
\end{proof}

\begin{theorem}[Exact neutral and wall nullities]
\label{cusp-native-full-flux-suN-nullity-theorem}
Away from an odd wall, the neutral Coulomb kernel has real dimension
\begin{equation}
 \operatorname{null}_{\mathrm{neutral}}\mathcal J_{\mathbf m,T}
 =4\left(\sum_{a=1}^{\ell}r_a^2-1\right).
 \label{cusp-native-full-flux-suN-neutral-nullity}
\end{equation}
At a wall \(\rho_T^*=2n+1>1\), the nonzero-root part adds
\begin{equation}
 12\sum_{i<j}q_{ij}^2=12r s_2(\mathbf m)
 \label{cusp-native-full-flux-suN-odd-wall-nullity}
\end{equation}
real zero modes.  At the equality wall \(S_T=0\), equivalently
\(\rho_T^*=1\), both favorable plane polarizations vanish and the nonzero
root contribution doubles to
\begin{equation}
 24\sum_{i<j}q_{ij}^2=24r s_2(\mathbf m).
 \label{cusp-native-full-flux-suN-equality-wall-nullity}
\end{equation}
The total nullity at a wall is the sum of the neutral term and the stated
root term.
\end{theorem}

\begin{proof}
On \(\mathfrak z_{\mathbf m}\), both the induced connection commutator and
the spin endomorphism vanish.  The operator is therefore the ordinary
coefficientwise Laplacian on one-forms.  Its kernel on a flat four-torus is
the space of constant one-forms, all of which are Coulomb.  Tensoring the
four-dimensional real harmonic-one-form space with
\(\mathfrak z_{\mathbf m}\) proves
\cref{cusp-native-full-flux-suN-neutral-nullity}.

If \(\rho_T^*=2n+1>1\), the high-plane favorable polarization at high
oscillator number zero and low oscillator number \(n\) is the unique zero
branch for each nonzero root.  Its positive-root-pair real dimension is
\(12q_{ij}^2\).  If \(S_T=0\), the two frequencies agree and the lowest
favorable polarization in each of the two planes is zero, so this dimension
doubles.  These charged zero modes are Coulomb: the intertwining identity
shows that their divergences lie in the kernel of the charged scalar
Laplacian, but that scalar Laplacian has strictly positive ground energy by
\cref{cusp-native-full-flux-suN-charge-d-levels}.  Summing over positive
roots and using
\(\sum_{i<j}q_{ij}^2=rs_2(\mathbf m)\) proves the result.
\end{proof}

For \(n\ge1\), the \(n\)-th odd wall is equivalently
\begin{equation}
 \rho_T^*=2n+1
 \iff n\Xi_T=(n+1)|S_T|
 \iff 24n^2D_T=(2n+1)S_T^2.
 \label{cusp-native-full-flux-suN-wall-equations}
\end{equation}
The final polynomial equation must be accompanied by the unsigned ordering
in the middle expression if one needs to distinguish the two sides
\(S_T>0\) and \(S_T<0\).  The \(n=0\) wall is exactly \(S_T=0\).

On the sufficiently deep cusp tail,
\(S_T>0\), \(5<\rho_T^*<7\), and \(N_-(T)=3\).  Hence
\begin{equation}
 \operatorname{ind}_{\mathbb R}\mathcal J_{\mathbf m,T}
 =36r s_2(\mathbf m)=6r k_{\mathbf m},
 \label{cusp-native-full-flux-suN-deep-index}
\end{equation}
and the three negative branches have the retained asymptotics
\begin{equation}
 \nu_{ij,n}(T)
 =\frac{q_{ij}\pi}{T}(4n-10)+O(T^{-2}),
 \qquad n=0,1,2.
 \label{cusp-native-full-flux-suN-negative-asymptotics}
\end{equation}
For example, the genuinely \(SU(3)\) data are
\begin{center}
\renewcommand{\arraystretch}{1.2}
\begin{tabular}{@{}lcccc@{}}
\toprule
weights & nonzero \(q_{ij}\) & \(k_{\mathbf m}\) & deep index
& neutral nullity\\
\midrule
\((1,0,-1)\) & \(1,1,2\) & \(12\) & \(216\) & \(8\)\\
\((1,1,-2)\) & \(3,3\) & \(36\) & \(648\) & \(16\)\\
\((2,1,-3)\) & \(1,4,5\) & \(84\) & \(1512\) & \(8\)\\
\bottomrule
\end{tabular}
\end{center}
The middle row exhibits both required edge phenomena: one neutral root from
the repeated weight and coincident numerical bands from the two charge-three
positive roots.  The specialization \(r=2\), \(\mathbf m=(1,-1)\) is
exactly the already proved native \(SU(2)\) theorem and is not reproduced
here.

\subsection{Stability boundary, possible wall bifurcation, and exact scope}

\begin{proposition}[Where the split background is, and is not, unstable]
\label{cusp-native-full-flux-suN-stability-boundary}
If \(S_T\ne0\), then \(\rho_T^*>1\), \(N_-(T)\ge1\), and every nonzero
positive root contributes a negative Coulomb direction.  Since a nonzero
trace-zero weight vector has at least one nonzero root,
\(\nabla^{\mathbf m,T}\) is then not a local minimum modulo gauge.  This
applies in particular throughout the sufficiently deep cusp tail.

If \(S_T=0\), the negative root spectrum is empty.  In that case
\cref{cusp-native-full-flux-suN-action-excess} attains the topological value
\(2\pi^2k_{\mathbf m}/g_{\mathrm{YM}}^2\), while
\cref{cusp-native-full-flux-suN-equality-wall-nullity} supplies root zero
modes.  No assertion of strict stability follows because the Hessian is
degenerate.
\end{proposition}

\begin{proof}
The first paragraph follows immediately from
\cref{cusp-native-full-flux-suN-ordered-frequencies,%
cusp-native-full-flux-suN-Morse-theorem}.  For the second,
\(S_T=0\) makes the two frequencies equal and hence \(\rho_T^*=1\); the
strict negative-band inequality has no solution.  The action and nullity
statements are exactly
\cref{cusp-native-full-flux-suN-action-excess,%
cusp-native-full-flux-suN-nullity-theorem}.
\end{proof}

A negative eigenspace is a space of descent directions, not a
Lyapunov--Schmidt bifurcation kernel.  A local bifurcation analysis must
instead be posed at one of the zero-mode parameters in
\cref{cusp-native-full-flux-suN-nullity-theorem}.  It must first quotient or
slice the ever-present centralizer-valued harmonic kernel, then identify the
remaining kernel and cokernel of the gauge-fixed nonlinear Yang--Mills map,
and finally compute the parameter transversality and nonlinear obstruction
coefficients.  The present linear calculation supplies the exact candidate
kernel at each wall; it does not assert existence of a nearby irreducible
branch.

The theorem proved here is strictly lower-rung and fibrewise.  It constructs
explicit reducible split-Cartan \(SU(r)\) Yang--Mills critical points on the
already existing regular flat four-torus fibres and determines their
topology, action, normalized Jacobi spectrum, physical Hessian scaling,
Morse index, and nullity.  It constructs no new base, spacetime, cone,
\(S^{12}\), \(S^{24}\), transfer matrix, or interacting quantum Hamiltonian.
The radial multiplication operator already treated in the canonical native
module is not repeated here.

There is also an exact obstruction to interpreting the nonzero-flux bundle
as the restriction of a bundle on a proposed higher carrier.  If
\(i:F_T\to Y\) is any map and \(H^4(Y;\mathbb Z)=0\), then every complex
bundle \(\mathscr V\to Y\) has \(c_2(\mathscr V)=0\), so naturality gives
\[
 c_2(i^*\mathscr V)=i^*c_2(\mathscr V)=0.
\]
This contradicts
\cref{cusp-native-full-flux-suN-Chern-classes,%
cusp-native-full-flux-suN-Chern-numbers} for
\(\mathscr E_{\mathbf m,T}\).  Therefore this nonzero-\(c_2\) bundle cannot
extend over any carrier with vanishing fourth integral cohomology, including
a space diffeomorphic to \(S^6\), \(S^{12}\), or \(S^{24}\).  Changing the
gauge group changes the root system on the fibre; it does not by itself
produce the proposed higher rung or any conclusion about a physical mass
gap.
